\section{Trajectory and Reward Computation}
\label{app:trajectory}

This appendix details the full trajectory capture and reward computation
pipeline that connects debate execution to GRPO training signal.

\subsection{DebateTrajectory Structure}

Each debate produces a \texttt{DebateTrajectory} that records every
decision and its outcome:

\begin{table}[h]
\centering
\small
\begin{tabular}{lp{8cm}}
\toprule
\textbf{Field} & \textbf{Description} \\
\midrule
\texttt{speeches} & Ordered list of all speeches (AC, AC-CX, NC, NC-CX, 1AR, NR, 2AR) with full text and metadata \\
\texttt{aff\_turns} & List of \texttt{TurnRecord}s for the affirmative side \\
\texttt{neg\_turns} & List of \texttt{TurnRecord}s for the negative side \\
\texttt{judge\_scores} & Per-speech scores from all sub-judges (Appendix~\ref{app:agentic_judge}) \\
\texttt{branches} & List of \texttt{BranchState} records (Appendix~\ref{app:branching}) \\
\texttt{outcome} & Final debate result: winner, score differential, and judge reasoning \\
\texttt{debate\_id} & Unique identifier linking to the resolution and side assignments \\
\bottomrule
\end{tabular}
\caption{DebateTrajectory top-level fields.}
\label{tab:trajectory}
\end{table}

\subsection{TurnRecord}

Each turn within the trajectory captures the full decision context:

\begin{table}[h]
\centering
\small
\begin{tabular}{lp{8cm}}
\toprule
\textbf{Field} & \textbf{Description} \\
\midrule
\texttt{tactic\_selection} & The macro-tactic and micro-tactic sequence chosen for this turn \\
\texttt{opponent\_prediction} & The model's predicted opponent response (used for strategic planning) \\
\texttt{opponent\_ground\_truth} & The actual opponent response observed after the turn \\
\texttt{novelty\_bonus} & Reward bonus for using a novel tactic combination ($+0.05$ to $+0.15$) \\
\texttt{surprise\_bonus} & Reward bonus when the opponent's response differs significantly from prediction ($+0.05$ to $+0.10$), indicating the model explored territory not covered by the opponent model \\
\texttt{call\_records} & Per-call records for tactic, skeleton, evidence, and speech generation, each with logprobs \\
\bottomrule
\end{tabular}
\caption{TurnRecord fields capturing per-turn decision context.}
\label{tab:turn_record}
\end{table}

The \texttt{opponent\_prediction} and \texttt{opponent\_ground\_truth}
fields enable the surprise bonus: when the model's prediction diverges
substantially from reality, it indicates unexplored strategic territory
that may yield high-signal training examples.

\subsection{Reward Formula}

The per-turn reward combines four components reflecting different aspects
of debate quality:

\begin{equation}
r_{\text{combined}} = 0.25 \cdot r_{\text{selection}} + 0.35 \cdot r_{\text{execution}} + 0.25 \cdot r_{\text{quality}} + 0.15 \cdot r_{\text{novelty}}
\label{eq:combined_reward}
\end{equation}

\begin{table}[h]
\centering
\small
\begin{tabular}{lcp{6.5cm}}
\toprule
\textbf{Component} & \textbf{Weight} & \textbf{Source} \\
\midrule
$r_{\text{selection}}$ & 0.25 & Tactic appropriateness for the debate state; from Forensics module (Appendix~\ref{app:agentic_judge}) \\
$r_{\text{execution}}$ & 0.35 & Speech quality conditional on the tactic; composite of sub-judge scores \\
$r_{\text{quality}}$ & 0.25 & Speech-specific primary dimension score (Table~\ref{tab:speech_dimensions}) \\
$r_{\text{novelty}}$ & 0.15 & Novelty bonus $+$ surprise bonus from \texttt{TurnRecord} \\
\bottomrule
\end{tabular}
\caption{Per-turn reward components. Execution receives the highest weight
because it reflects the most direct measure of output quality.}
\label{tab:reward_components}
\end{table}

The final reward incorporates the debate outcome:
\begin{equation}
r_{\text{final}} = 0.7 \cdot r_{\text{combined}} + 0.3 \cdot r_{\text{win}}
\label{eq:final_reward}
\end{equation}

\noindent where $r_{\text{win}} \in \{0, 1\}$ is 1 if the model's side
won the debate and 0 otherwise. The 0.3 weight on the win bonus ensures
that outcome signal influences training without overwhelming the
fine-grained per-turn signal. Without the win bonus, the model could
learn to produce individually high-scoring speeches that collectively
lose the debate (e.g., strong individual arguments that fail to
cohere into a winning narrative).

\subsection{BackpropScorer Algorithm}

The BackpropScorer propagates debate-level outcomes back to individual
calls, producing causal attribution of which decisions most influenced
the result. This operates in four stages:

\paragraph{Stage 1: Outcome analysis.}
For each speech in the debate, compute its \emph{outcome contribution}
$c_s \in [-1.0, +1.0]$, estimating how much that speech contributed to
the final result. This is assessed by the scoring LLM, which receives
the full debate transcript and is asked: ``How much did [speech] contribute
to [side] winning/losing? Consider what would have changed if this speech
had been average rather than as-delivered.''

\paragraph{Stage 2: Branch counterfactual.}
For debates with branches (Appendix~\ref{app:branching}), identify the
divergence point and determine whether a winner shift occurred. If so,
the speeches at the divergence point receive amplified attribution:
their contribution scores are adjusted by the score differential between
the winning and losing branches.

\paragraph{Stage 3: Score adjustment.}
The final backpropagated score for each call is:
\begin{equation}
r_{\text{backprop}} = r_{\text{inline}} + c_s \cdot 0.2 + b_{\text{critical}} \cdot 0.15 + p_{\text{adj}}
\label{eq:backprop}
\end{equation}

\begin{table}[h]
\centering
\small
\begin{tabular}{lp{8cm}}
\toprule
\textbf{Term} & \textbf{Description} \\
\midrule
$r_{\text{inline}}$ & Original inline score from the sub-judges \\
$c_s \cdot 0.2$ & Outcome contribution scaled by 0.2 to prevent outcome signal from dominating \\
$b_{\text{critical}} \cdot 0.15$ & Critical call boost: $b_{\text{critical}} \in \{0, 1\}$ is 1 if Forensics identifies this call as a pivotal decision point \\
$p_{\text{adj}}$ & Pivotal adjustment from branch counterfactuals: ranges from $-0.5$ to $+0.5$ for winner-shifted decisions, 0 otherwise \\
\bottomrule
\end{tabular}
\caption{BackpropScorer adjustment terms.}
\label{tab:backprop_terms}
\end{table}

The multiplicative constants (0.2, 0.15) were tuned to keep
$r_{\text{backprop}}$ within approximately $\pm 0.3$ of
$r_{\text{inline}}$ for most calls, with larger adjustments reserved
for the small number of genuinely pivotal decisions (typically 2--3
per debate).

\paragraph{Stage 4: Reasoning preservation.}
The BackpropScorer generates a natural-language explanation for each
score adjustment, recording the causal chain from the call to the
debate outcome. These explanations are stored alongside the training
data for potential use in reasoning distillation---training the model
to internalize the judge's causal reasoning rather than just the
scalar reward.

\subsection{Group Baselines and Context Triplets}

As described in Appendix~\ref{app:grpo_groups}, GRPO advantages are
normalized within context triplets
$(\text{our\_tactic}, \text{opponent\_tactic}, \text{opponent\_intent})$.
The trajectory captures these triplets at each turn, enabling the
training pipeline to compute context-appropriate baselines.

\begin{table}[h]
\centering
\small
\begin{tabular}{lp{5cm}r}
\toprule
\textbf{Triplet Component} & \textbf{Source} & \textbf{Cardinality} \\
\midrule
Our tactic & \texttt{tactic\_selection} from TurnRecord & $\sim$25 macro-tactics \\
Opponent tactic & Inferred from opponent speech by the judge & $\sim$25 macro-tactics \\
Opponent intent & Classified as \textsc{Attack}, \textsc{Defend}, \textsc{Extend}, or \textsc{Probe} & 4 \\
\bottomrule
\end{tabular}
\caption{Context triplet components for advantage normalization. The
product space is large ($\sim$2,500 triplets) but most debates activate
only 15--30 distinct triplets, providing sufficient samples for
meaningful normalization.}
\label{tab:context_triplets}
\end{table}

\subsection{End-to-End Reward Flow}

Figure~\ref{fig:reward_flow} summarizes the complete reward computation
pipeline from debate execution to GRPO training signal.

\begin{figure}[h]
\centering
\small
\begin{tabular}{c}
\hline
\\[-0.5em]
\textbf{Debate Execution} \\
$\downarrow$ \\
\texttt{DebateTrajectory} (speeches, turns, branches) \\
$\downarrow$ \\
\textbf{Layer 1: Inline Scoring} \\
4 sub-judges $\times$ speech-specific dimensions $\rightarrow$ $r_{\text{inline}}$ \\
$\downarrow$ \\
\textbf{Layer 2: Outcome Backpropagation} \\
$r_{\text{inline}} + \text{outcome contribution} + \text{critical boost} + \text{branch adjustment}$ $\rightarrow$ $r_{\text{backprop}}$ \\
$\downarrow$ \\
\textbf{Layer 3: Meta-Debate Calibration} \\
3-model panel $+$ bias detection $+$ judge-the-judge filtering \\
$\downarrow$ \\
\textbf{Reward Assembly} \\
$r_{\text{final}} = 0.7 \cdot r_{\text{combined}} + 0.3 \cdot r_{\text{win}}$ \\
$\downarrow$ \\
\textbf{Triplet Normalization} \\
$A_i = r_i - \bar{r}_{(\text{our\_tactic}, \text{opp\_tactic}, \text{opp\_intent})}$ \\
$\downarrow$ \\
\textbf{GRPO Training} (Group A $\rightarrow$ B $\rightarrow$ C $\rightarrow$ D) \\
\\[-0.5em]
\hline
\end{tabular}
\caption{End-to-end reward flow from debate execution to GRPO training.
Each layer adds finer-grained signal: inline scoring evaluates individual
speeches, backpropagation attributes outcomes to decisions, and
meta-debate calibration removes judge bias.}
\label{fig:reward_flow}
\end{figure}
